\subsubsection*{Solution}
\paragraph*{Interpret the requested ordering literally}

\section{Conventions and plan}
Throughout, $d=8$ at the end but $d$ is kept general where possible, $L$ is the AdS radius (set to $1$ inside curvature polynomials), curvature conventions are standard (round spheres positive), and the boundary metric is Euclidean. I write
\begin{equation}
S_{\mu\nu}=\frac{1}{d-2}\Big(R_{\mu\nu}-\frac{R}{2(d-1)}\gamma_{\mu\nu}\Big),\qquad J=S^\mu{}_\mu ,
\end{equation}
for the usual Schouten tensor, so that the problem's $P_{\mu\nu}=(d-2)S_{\mu\nu}$ ($=6S_{\mu\nu}$ in $d=8$). With this, the problem's $B_{\mu\nu}$ is the Bach tensor $B_{\mu\nu}=\nabla^\rho C^S_{\mu\nu\rho}-W_{\rho\nu\mu\sigma}S^{\sigma\rho}$ and the problem's $O_{\mu\nu}$ is the six-dimensional obstruction tensor $O^{(6)}_{\mu\nu}$, related to Graham's extended obstruction tensor by $O_{\mu\nu}=(d-4)(d-6)\Omega^{(2)}_{\mu\nu}$; $\Omega^{(1)}_{\mu\nu}=-\frac{1}{d-4}B_{\mu\nu}$. The tensor $\Omega_{\mu\nu}$ of the problem is taken literally as defined.

Plan: (1) show that the anomaly is the coefficient of $\rho^{d/2}$ in $\sqrt{\det g(x,\rho)}$ in Fefferman--Graham (FG) gauge, and fix the normalisation and sign of $X^{(4)}$ from eqs.~(1) and (8) of the problem; (2) obtain the FG coefficients $g_{(2)},g_{(4)},g_{(6)}$; (3) use the $\rho\rho$ Einstein equation to write the $\rho^4$ coefficient of $\sqrt{\det g}$ in terms of traces; (4) assemble $X^{(4)}$ and rewrite it in the problem's variables; (5) verify everything by an exact computer calculation.

\section{The anomaly as the logarithmic term}
Write the bulk metric in FG gauge,
\begin{equation}
ds^2=L^2\Big(\frac{d\rho^2}{4\rho^2}+\frac1\rho g_{ij}(x,\rho)dx^idx^j\Big),\qquad
g(x,\rho)=g_{(0)}+\rho g_{(2)}+\rho^2g_{(4)}+\rho^3g_{(6)}+\cdots,
\end{equation}
and define $v(x,\rho)=\sqrt{\det g(x,\rho)/\det g_{(0)}}=1+\rho v^{(2)}+\rho^2v^{(4)}+\cdots$. On shell, $R=-d(d+1)/L^2$, so the Euclidean action $S=-\frac{1}{16\pi G}\int\sqrt G\,(R+\frac{d(d-1)}{L^2})+S_{\rm GH}$ becomes
\begin{equation}
S_{\rm bulk}=\frac{dL^{d-1}}{16\pi G}\int d^dx\sqrt{g_{(0)}}\int_\epsilon d\rho\,\rho^{-d/2-1}v(x,\rho)\ \ni\ -\frac{dL^{d-1}}{16\pi G}\ln\epsilon\int d^dx\sqrt{g_{(0)}}\,v^{(d)} .
\end{equation}
The Gibbons--Hawking term and the local counterterms contribute no $\ln\epsilon$ (the $\rho^{d/2}\ln\rho$ term of $g$ is traceless and does not enter $v$ at this order). A constant Weyl rescaling $g_{(0)}\to e^{2\sigma}g_{(0)}$ is the same bulk solution described with $\rho\to e^{2\sigma}\rho$, i.e.\ with the cutoff $\epsilon\to e^{-2\sigma}\epsilon$; the renormalised action, which contains the explicit log counterterm $+\frac{dL^{d-1}}{16\pi G}\ln\epsilon\int\sqrt{\gamma}\,v^{(d)}[\gamma]$, therefore changes by
\begin{equation}
S_{\rm ren}[e^{2\sigma}g_{(0)}]-S_{\rm ren}[g_{(0)}]=\frac{2\sigma\,dL^{d-1}}{16\pi G}\int d^dx\sqrt{g_{(0)}}\,v^{(d)} .
\end{equation}
With $Z=e^{-S_{\rm ren}}$ and the convention adopted for eq.~(1) of the problem, $Z[\gamma]=e^{-\mathcal A}Z[\mathcal B^{-2}\gamma]$, i.e.\ $\mathcal A=S_{\rm ren}[\gamma]-S_{\rm ren}[\mathcal B^{-2}\gamma]$; taking $e^{2\sigma}=\mathcal B^{-2}$, $\sigma=-\ln\mathcal B$ (and extending to local $\mathcal B$, since $\sqrt\gamma v^{(d)}$ is the local anomaly density),
\begin{equation}
\mathcal A=\frac{dL^{d-1}}{8\pi G}\int d^dx\sqrt{\gamma}\;v^{(d)}\ln\mathcal B .
\end{equation}
Comparing with eq.~(8) of the problem, $\mathcal A_4=-\frac{L^7}{8\pi G}\int\sqrt\gamma\,X^{(4)}\ln\mathcal B$,
\begin{equation}
\boxed{X^{(k)}=-2k\,v^{(2k)}\big|_{d=2k}},\qquad X^{(4)}=-8\,v^{(8)}\big|_{d=8}.
\label{eq:X}
\end{equation}
\emph{Checks.} $d=2$: $v^{(2)}=\frac12\tr g_{(2)}=-\frac12J=-\frac{R}{4}$, so $X^{(1)}=\frac R2$ and $\mathcal A_1=-\frac{L}{16\pi G}\int\sqrt\gamma R\ln\mathcal B$, i.e.\ $\langle T^\mu{}_\mu\rangle=-\frac{c}{24\pi}R$ with $c=\frac{3L}{2G}$, the convention used in the source. $d=4$: $v^{(4)}=-\frac18\tr g_{(2)}^2+\frac18(\tr g_{(2)})^2=-\frac1{32}\big(R_{\mu\nu}R^{\mu\nu}-\frac13R^2\big)$, so $X^{(2)}=\frac18(R_{\mu\nu}R^{\mu\nu}-\frac13R^2)$, the Henningson--Skenderis result with $a=c$. (With the other reading of eq.~(1) all signs flip.)

\section{Fefferman--Graham coefficients}
The $ij$ components of the Einstein equations in FG gauge read (standard curvature signs; primes are $\partial_\rho$)
\begin{equation}
\rho\big[2g''-2g'g^{-1}g'+\tr(g^{-1}g')\,g'\big]-{\rm Ric}(g)-(d-2)g'-\tr(g^{-1}g')\,g=0,
\label{eq:FG}
\end{equation}
which I checked against the exact solution $g(\rho)=(1-\rho/4)^2g_{S^d}$ of hyperbolic space. At order $\rho^0$: $g_{(2)}=-S$. At order $\rho^1$ one needs the linearised Ricci tensor; using
$\nabla^k\nabla_iS_{jk}=\nabla_i\nabla_jJ-W_{kimj}S^{km}+dS^2_{ij}-|S|^2g_{ij}$ (commute derivatives, Bianchi identity, Weyl decomposition) one finds
\begin{equation}
g_{(4)}=\tfrac14\Big(S^2-\frac{1}{d-4}B\Big),\qquad B_{ij}=\nabla^kC_{ijk}-W_{kijl}S^{kl}=\Box S_{ij}-\nabla_i\nabla_jJ-2W_{kijl}S^{kl}-dS^2_{ij}+|S|^2g_{ij},
\end{equation}
in agreement with Fefferman--Graham. At order $\rho^2$ the same procedure gives $g_{(6)}$; the result is most compactly obtained from the ambient-metric expansion (Challenge 53): with $\rho_{\rm amb}=-\rho/2$, $g_{(6)}=-\frac18\gamma^{(3)}$ and $\gamma^{(3)}=\frac13\Omega^{(2)}+\frac{4}{3(4-d)}B_{k(i}S^k{}_{j)}$, so
\begin{equation}
g_{(6)}=-\frac{1}{24}\Omega^{(2)}+\frac{1}{6(d-4)}(BS)_{(ij)}
=-\frac{1}{24(d-6)(d-4)}O_{ij}+\frac{1}{6(d-4)}B_{k(i}S^k{}_{j)} .
\label{eq:g6}
\end{equation}
This identifies the problem's $O_{\mu\nu}$ (eq.~(4) of the problem, written with $P=(d-2)S$) with $(d-4)(d-6)\Omega^{(2)}_{\mu\nu}$; the identification, and (\ref{eq:g6}) as an exact tensor identity in $d=8$, were verified by the computer calculation of Section~\ref{sec:num}.

\section{$v^{(8)}$ from the $\rho\rho$ equation}
The $\rho\rho$ Einstein equation is $\tr(g^{-1}g'')=\frac12\tr\big((g^{-1}g')^2\big)$. Since $(\ln\det g)''=\tr(g^{-1}g'')-\tr((g^{-1}g')^2)$, this is equivalent to
\begin{equation}
(\ln v)''=-\tfrac14\tr\big((g^{-1}g')^2\big),\qquad (\ln v)'(0)=\tfrac12\tr g_{(2)} .
\end{equation}
Writing (matrix notation, indices raised with $g_{(0)}$) $g_1=g_{(2)}$, $g_2=g_{(4)}$, $g_3=g_{(6)}$, one has $g^{-1}g'=g_1+\rho(2g_2-g_1^2)+\rho^2(3g_3-2g_1g_2-g_2g_1+g_1^3)+\dots$, hence
\begin{equation}
\tr\big((g^{-1}g')^2\big)=\tr g_1^2+2\rho\,\tr(2g_1g_2-g_1^3)+\rho^2\big[4\tr g_2^2-10\tr(g_1^2g_2)+3\tr g_1^4+6\tr(g_1g_3)\big]+\dots
\end{equation}
Integrating twice, $\ln v=\sum_k c_k\rho^k$ with
\begin{equation}
c_1=\tfrac12\tr g_1,\quad c_2=-\tfrac18\tr g_1^2,\quad c_3=-\tfrac16\tr(g_1g_2)+\tfrac1{12}\tr g_1^3,\quad
c_4=-\tfrac1{12}\tr g_2^2+\tfrac5{24}\tr(g_1^2g_2)-\tfrac1{16}\tr g_1^4-\tfrac18\tr(g_1g_3),
\end{equation}
and $v^{(8)}=c_4+c_1c_3+\frac12c_2^2+\frac12c_1^2c_2+\frac1{24}c_1^4$. (Consistency: this reproduces $\tr g_{(4)}=\frac14\tr g_{(2)}^2$ and $\tr g_{(6)}=\frac{1}{6(d-4)}\tr(BS)$, which also follow from the trace of (\ref{eq:FG}) and are trace-free conditions on $B$ and $\Omega^{(2)}$.)

\section{Assembling $X^{(4)}$}
Insert $g_1=-S$, $g_2=\frac14(S^2-\frac{B}{d-4})$ and $\tr(g_1g_3)=\frac1{24}\tr(S\Omega^{(2)})-\frac{1}{6(d-4)}\tr(BS^2)$ from (\ref{eq:g6}) (using $\tr(S\,(BS)_{\rm sym})=\tr(BS^2)$):
\begin{align}
c_1&=-\tfrac J2,\qquad c_2=-\tfrac18\tr S^2,\qquad c_3=-\tfrac1{24}\tr S^3-\tfrac{1}{24(d-4)}\tr(BS),\\
c_4&=-\tfrac1{64}\tr S^4-\tfrac{1}{48(d-4)}\tr(BS^2)-\tfrac{1}{192(d-4)^2}\tr B^2-\tfrac1{192}\tr(S\Omega^{(2)}),
\end{align}
so that
\begin{equation}
\begin{split}
v^{(8)}=&-\tfrac1{64}\tr S^4+\tfrac{J}{48}\tr S^3+\tfrac1{128}(\tr S^2)^2-\tfrac{1}{64}J^2\tr S^2+\tfrac1{384}J^4
+\tfrac{J}{48(d-4)}\tr(BS)\\
&-\tfrac{1}{48(d-4)}\tr(BS^2)-\tfrac{1}{192(d-4)^2}\tr B^2-\tfrac1{192}\tr(S\Omega^{(2)}).
\end{split}
\end{equation}
Setting $d=8$ and using (\ref{eq:X}):
\begin{equation}
X^{(4)}=\tfrac18\tr S^4-\tfrac16 J\tr S^3-\tfrac1{16}(\tr S^2)^2+\tfrac18J^2\tr S^2-\tfrac1{48}J^4
-\tfrac1{24}J\tr(BS)+\tfrac1{24}\tr(BS^2)+\tfrac1{384}\tr B^2+\tfrac1{24}\tr(S\,\Omega^{(2)}) .
\label{eq:XS}
\end{equation}
This is an exact local identity in FG gauge (no total derivatives are dropped), and it coincides with eq.~(72) of Jia--Karydas once $\Omega^{(1)}=-\frac14B$ is inserted there. Finally translate to the problem's variables, $S=P/6$, $J=\frac16\tr P$, $\Omega^{(2)}=\frac18O$ ($d=8$):
\begin{align}
X^{(4)}=\;&\frac{1}{10368}\tr(P^4)-\frac{1}{7776}\tr(P^3)\tr(P)-\frac{1}{20736}\big(\tr P^2\big)^2+\frac{1}{10368}\tr(P^2)\big(\tr P\big)^2-\frac{1}{62208}\big(\tr P\big)^4\nonumber\\
&-\frac{1}{864}\tr(BP)\tr(P)+\frac{1}{864}\tr(BP^2)+\frac{1}{384}\tr(B^2)+\frac{1}{1152}\tr(OP).
\label{eq:XP}
\end{align}
All other candidate terms of the problem have coefficient zero: $\tr P^3$, $\tr BP$, $\tr\Omega$ (weight 6) and $\tr B^2P$, $\tr OP^2$ (weight 10) cannot appear in a weight-8 density, and $\tr(\Omega P)$, although of weight 8, does not appear: $\Omega_{\mu\nu}$ contains the structures $\Box B$, $WB$, $BJ$, $P\nabla C$, $CC$, \dots with relative coefficients different from those in $O_{\mu\nu}$, so no combination $\alpha\,O+\beta\,\Omega$ with $\beta\neq0$ can reproduce the $\Box B$ and $\nabla C$ structures of (\ref{eq:XS}) simultaneously (the ratio of their coefficients differs by a factor $36$ between $O$ and $\Omega$), and the remaining listed terms contain no such structures. This is confirmed below by the exact linear algebra.

\section{Exact numerical verification}\label{sec:num}
To guard against sign and factor errors in the tensor algebra, I verified (\ref{eq:XP}) by an exact computation:
\begin{itemize}
\item $23$ random metrics $g_{(0)ij}(x)=\delta_{ij}+(\text{random integer polynomial of degree}\le6)$ in $d=8$ were generated; all tensor algebra was carried out with truncated multivariate Taylor polynomials over $\mathbb F_p$, $p=2^{31}-1$ (exact arithmetic, no rounding).
\item The Einstein equations (\ref{eq:FG}) were solved order by order for $g_{(2)},g_{(4)},g_{(6)}$ directly (Ricci tensors of $g(x,\rho)$ computed from scratch). For every metric the following identities held exactly: $g_{(2)}=-S$; $\tr g_{(4)}=\frac14\tr g_{(2)}^2$; $g_{(4)}=\frac14(S^2-\frac14B)$ with $B$ computed from the problem's eq.~(3); $\tr g_{(6)}=\frac1{24}\tr(BS)$; $g_{(6)}=-\frac{1}{192}O+\frac1{24}B_{k(i}S^k{}_{j)}$ with $O$ computed from the problem's eqs.~(4)--(5); $\tr O=0$.
\item $X^{(4)}=-8v^{(8)}$ was computed from the traces, and all seventeen invariants ($\tr P^4$, $\tr P^3$, $\tr P^3\tr P$, $\tr BP$, $\tr BP^2$, $\tr B^2$, $\tr B^2P$, $\tr OP$, $\tr OP^2$, $\tr\Omega$, $\tr\Omega P$, $(\tr P^2)^2$, $\tr P^2(\tr P)^2$, $(\tr P)^4$, $\tr(BP)\tr P$, $\tr P^2\tr P$, $(\tr P)^3$) were evaluated at the origin; $\Omega$ was computed from the problem's eqs.~(6)--(7).
\item The linear system for the coefficients of the nine weight-8 invariants has rank $9$ and is consistent over all $23$ samples; the rational reconstruction of its unique solution is exactly (\ref{eq:XP}). Adding $\tr(\Omega P)$ as a tenth unknown keeps the system of full rank (so the ten functions are linearly independent) and returns coefficient $0$ for it. Restricting to the eleven printed terms gives an \emph{inconsistent} system (rank 10), confirming that the printed list is incomplete. Incidentally, $\tr\Omega=\tr P^3$ held identically for all samples (the derivative terms in $\tr\Omega$ cancel by $C^{\rho\mu\lambda}C_{\lambda\mu\rho}=\nabla^\lambda P^{\rho\mu}C_{\mu\rho\lambda}$), which is the only linear relation among the seventeen invariants.
\end{itemize}

\section{Result}
In the order printed in the problem,
\begin{equation}
\begin{split}
&\big(\tr P^4,\ \tr P^3,\ \tr P^3\tr P,\ \tr BP,\ \tr BP^2,\ \tr B^2,\ \tr B^2P,\ \tr OP,\ \tr OP^2,\ \tr\Omega,\ \tr\Omega P\big)\\
&\mapsto\
\Big(\tfrac{1}{10368},\,0,\,-\tfrac{1}{7776},\,0,\,\tfrac{1}{864},\,\tfrac{1}{384},\,0,\,\tfrac{1}{1152},\,0,\,0,\,0\Big),
\end{split}
\end{equation}
together with the four additional terms of (\ref{eq:XP}): $(\tr P^2)^2:-\frac{1}{20736}$, $\ \tr P^2(\tr P)^2:\frac{1}{10368}$, $\ (\tr P)^4:-\frac{1}{62208}$, $\ \tr(BP)\tr P:-\frac{1}{864}$.

\section{Comparison with the provided solution and corrections}
The provided solution quotes the eight-dimensional formula of Jia--Karydas in conventional Schouten normalisation and converts it; its eleven coefficients agree with mine, and its observation that the printed list is incomplete is correct. \chg{The following points are corrected or completed here:
\begin{enumerate}
\item The provided solution does not derive the formula it quotes, nor does it fix the overall sign from the problem's own definitions. Sections~2--5 supply a first-principles derivation; in particular, the sign is shown to follow from reading eq.~(1) of the problem as $Z[\gamma]=e^{-\mathcal A}Z[\mathcal B^{-2}\gamma]$ (the source's convention), and the reader is warned that the opposite reading flips all signs.
\item The provided solution asserts the relations $\Omega^{(1)}=-\frac14B$, $\Omega^{(2)}=\frac18O$ ``from Graham''. These are correct, but only because the problem's $O_{\mu\nu}$ (with its modified $\frac{1}{d-2}$ factors) is exactly $O^{(6)}$; this is established here analytically (through $g_{(6)}$) and verified exactly by computer algebra, which also confirms the problem's $B_{\mu\nu}$ (eq.~(3)) as the Bach tensor.
\item The provided solution silently ignores the problem's $\Omega_{\mu\nu}$. Here it is pointed out that $\Omega_{\mu\nu}$ as printed is not the extended obstruction tensor, that $\tr\Omega=\tr P^3$ identically, and that the vanishing of the $\tr(\Omega P)$ coefficient is nonetheless unambiguous (full-rank linear independence).
\item The complete density (\ref{eq:XP}), including the four terms missing from the printed list, is given explicitly and verified.
\end{enumerate}
The final numerical values of the requested eleven coefficients are unchanged.}